\begin{lemma}
\label{editorial-source-lemma-flat-fibres-irreducible}
Let $R \to S$ be a ring map. Assume
\begin{enumerate}
\item[(a)] $\Spec(R)$ is irreducible,
\item[(b)] $R \to S$ is flat,
\item[(c)] $R \to S$ is of finite presentation,
\item[(d)] the fibre rings $S \otimes_R \kappa(\mathfrak p)$
have irreducible spectra for a dense collection of primes $\mathfrak p$ of $R$.
\end{enumerate}
Then $\Spec(S)$ is irreducible.
This is true more generally with (b) $+$ (c)
replaced by ``the map $\Spec(S) \to \Spec(R)$ is open''.
\end{lemma}

\begin{proof}
Write $\varphi:R\to S$ for the original
ring map. Put $X=\Spec(S)$, $Y=\Spec(R)$
and retain its spectrum map $f:X\to Y$. Proposition
\ref{editorial-source-proposition-fppf-open} makes $f$ open
under (b) and (c). We prove the stated
more general open-map assertion directly.
Let $D\subseteq Y$ be the given dense
set of points with irreducible fibres.
The irreducible space $Y$ is nonempty,
so $D$ is nonempty. Each fibre at a
point of $D$ is nonempty; consequently
$X$ is nonempty as well.

Let $U,V$ be nonempty open subsets of
$X$. The original images $f(U),f(V)$
are nonempty open subsets of the
irreducible $Y$, so their intersection
is nonempty and open. Density supplies
$y\in D\cap f(U)\cap f(V)$.
The subsets $U\cap f^{-1}(y)$ and
$V\cap f^{-1}(y)$ are nonempty open
subsets of the irreducible fibre.
Their intersection is nonempty,
so $U\cap V\neq\emptyset$.
This is the nonempty-space criterion
for irreducibility, proving the lemma
and supplying the nonemptiness step
in the cited Topology, Lemma
\ref{topology-lemma-irreducible-on-top}.

The same proof needs less than
openness: it is enough that the
image of every nonempty open
$U\subseteq X$ contain a nonempty
open subset of $Y$. Choose such
subsets inside $f(U)$ and $f(V)$
and perform the same intersection
argument with $D$. The whole proof
also works with the actual subspace
$f(X)\subseteq Y$ as codomain,
provided it is irreducible, the
specified interior condition is
relative to that subspace, and
the good fibres form a dense
subset of it. These are statements
about the original map with its
codomain restricted, not an
assumption that the image equals $Y$.

For flat ring maps there is a
different exact criterion without
finite presentation. Under (a),
let $\mathfrak p=\operatorname{Nil}(R)$
be the unique minimal prime of
the original $R$. Put
$C=S\otimes_R\kappa(\mathfrak p)$.
The original fibre comparison is
$$
C\simeq(R\setminus\mathfrak p)^{-1}(S/\mathfrak pS),
\qquad
s\otimes(\overline a/\overline b)
\longmapsto\varphi(a)s/\varphi(b),
\quad b\notin\mathfrak p.
$$
The inverse sends $\overline s/\varphi(b)$
to $s\otimes(1/\overline b)$;
quotient and localization balancing
prove well-definedness and both
composites. Primes of $C$ therefore
correspond exactly to original
primes of $S$ contracting to
$\mathfrak p$.

If $\mathfrak q$ is a minimal
prime of $S$, going down for
flat maps (Lemma \ref{editorial-source-lemma-flat-going-down}) lifts the inclusion
$\mathfrak p\subseteq\varphi^{-1}(\mathfrak q)$
to a prime contained in
$\mathfrak q$. Minimality forces
equality, so its contraction is
exactly $\mathfrak p$. Its fibre
prime is minimal, since a smaller
fibre prime would contract to
a smaller prime of $S$.
Conversely, if a minimal fibre
prime corresponds to $\mathfrak q$,
choose a minimal prime
$\mathfrak q_0\subseteq\mathfrak q$
of $S$. Its contraction is
$\mathfrak p$ by the previous
argument. Its fibre prime is
contained in the given minimal
one, so both are equal, and
$\mathfrak q_0=\mathfrak q$.
Thus these original fibre maps
give a bijection between the
minimal primes of $S$ and $C$.
If either ring is zero, both
sets are empty; if $S$ is nonzero,
it has a minimal prime and hence
$C$ is nonzero. Consequently,
under (a) and flatness alone,
$\Spec(S)$ is irreducible if and
only if this original generic
fibre has irreducible spectrum.

Dense irreducible fibres cannot
replace the generic-fibre condition
after dropping the openness
assumption. Retain the original
diagonal map
$$
\mathbf Z\longrightarrow\mathbf Q\times\mathbf Z,
\qquad n\longmapsto(n,n).
$$
It is flat: its underlying module
is the finite direct sum of the
flat localization $\mathbf Q$
and the free module $\mathbf Z$.
For every prime integer $\ell$,
its fibre is $\mathbf F_\ell$.
Indeed $\mathbf Q\otimes_{\mathbf Z}\mathbf F_\ell=0$,
since $1\otimes1=\ell^{-1}\otimes\ell=0$,
and the second factor gives
$\mathbf Z\otimes\mathbf F_\ell\simeq\mathbf F_\ell$.
These closed primes are dense:
every nonempty $D(n)$, $n\neq0$,
contains a prime dividing $|n|+1$
and hence not dividing $n$.
But the generic fibre is
$\mathbf Q\times\mathbf Q$,
through the multiplication map
$q\otimes(a,z)\mapsto(qa,qz)$
and its inverse
$(a,b)\mapsto a\otimes(1,0)+b\otimes(0,1)$.
The equality on the first factor
uses the original localization
identity $\mathbf Q\otimes_{\mathbf Z}\mathbf Q\simeq\mathbf Q$.
Both the total spectrum and its
generic fibre are reducible.
The nonempty open component
$\Spec(\mathbf Q)$ has image
$\{(0)\}\subseteq\Spec(\mathbf Z)$,
which has empty interior by the
same $D(n)$ argument. This
identifies exactly the failed
topological hypothesis.
\end{proof}